\documentclass[a4paper,11pt]{article}

\usepackage[german]{babel}
\usepackage[applemac]{inputenc}
\usepackage{listings}
\usepackage{graphicx}
\usepackage{multicol}
\title{Timespan Converter}
\begin{document}

\maketitle

\section{Einf�hrung}
Schreiben Sie ein Programm {\tt TimespanConverter}, in welchem Minuten in Wochen, Tage, Stunden und Minuten umgewandelt werden. Die Ausgabe soll so erfolgen, dass vorkommende Nullen nicht ausgegeben werden. Zur Verdeutlichung betrachten Sie folgende Beispiele:

\begin{center}
\includegraphics[scale=.7]{SampleOutput.png}
\end{center}

\section{Hinweise}
\begin{itemize}
	\item Wie im zweiten und im letzten Beispiel zu sehen ist, werden, falls die eingegebenen Minuten bei der Umwandlung 0 ergeben, diese nicht ausgegeben. 40320 Minuten sind exakt 4 Wochen, daher werden die Ausgabe von 0 Tage, 0 Stunden, 0 Minuten unterdr�ckt.
	
	\item Wie im zweiten Beispiel zu sehen ist, m�ssen Sie nicht auf die Ausgabe von Einheiten in der Einzahl achten. Auch wenn es korrekterweise {\tt 60 Minuten sind 1 Stunde} hei�en m�sste, k�nnen Sie hier das Wort {\em Stunde} in der Mehrzahl lassen. Gleiches gilt auch f�r Tage und Wochen.
	
	\item Fehlerhafte Eingaben, wie nicht ganzzahlige Minutenangaben oder negative Minutenangaben brauchen Sie nicht zu ber�cksichtigen.
	
	\item Das Beispiel ist mit Notepad++ und mit Hilfe des Compileraufrufs in der Eingabeaufforderung zu l�sen.
\end{itemize}

\section{L�sungsidee}
\begin{enumerate}
	\item Sie speichern die eingegebenen Minuten in eine neue Variable (beispielsweise {\tt workMinutes}).
	\item Sie dividieren {\tt workMinutes} durch 10 080 und erhalten die Wochen.
	\item Sie subtrahieren 10 080 von {\tt workMinutes}.
	\item Sie dividieren {\tt workMinutes}  nun durch 1 440 und erhalten die Tage.
	\item Sie subtrahieren 1 440 von {\tt workMinutes}.
	\item Sie dividieren {\tt workMinutes} durch 60 und erhalten die Stunden.
	\item Sie ermitteln den Rest der Division von {\tt workMinutes} durch 60 und erhalten die Minuten.
	\item Sie geben nun die in den Schritten 1 bis 7 erhaltenen Werte aus, allerdings nur dann, wenn der jeweilige Wert gr��er als 0 ist.
\end{enumerate}
\section{Source Code Qualit�t}
Achten Sie darauf, dass Sie korrekte Einr�ckungen machen und dass die Variablennamen richtig geschrieben werden (Gro�- und Kleinschreibung).

\end{document}